
%%%%%%%%%%%%%%%%%%%%%%%%%%%%
%%                        %%
%%         Resumo         %%
%%                        %%
%%%%%%%%%%%%%%%%%%%%%%%%%%%%

% Title format for the abstract: set time new roman as font of the index content


\renewenvironment{resumo}{\vspace{15pt}
    \normalsize
    {\normalfont\bfseries\Large\textbf{RESUMO}}
    \vspace{15pt}
}{\endquotation}

\renewenvironment{abstract}{\vspace{15pt}
    \normalsize\centering
    {\normalfont\bfseries\Large\textbf{ABSTRACT}}
    \vspace{15pt}
}{\endquotation}

\begin{resumo}
	\noindent

    \justifying

    Este trabalho tem como objetivo apresentar uma proposta de desenvolvimento de uma plataforma de análise e visualização de dados extensível e personalizável, visando proporcionar uma forma simples e objetiva de acesso para as bibliotecas de aprendizado de máquina e semelhantes presentes no repositório de bibliotecas \textit{Pypi}. Para atingir este objetivo, foi realizado entrevistas com os pesquisadores da área de Ciência de Dados e Processamento de Sinais, responsáveis pelo projeto de pesquisa \textit{SmartEnergy}, para identificar as principais necessidades e desafios no processo de análise de dados. Além disso, foi realizada uma pesquisa bibliográfica para identificar as principais ferramentas de análise de dados e visualização de dados que possam atender aos requisitos identificados. Por fim, foi proposto um conjunto de bibliotecas e ferramentas para o desenvolvimento da plataforma, que podem ser utilizadas para a criação de um ambiente de desenvolvimento para análise e visualização de dados, que atenda as necessidades do processo de análise de dados.

    \textbf{Palavras-chave}: Plataforma, Análise de Dados, Visualização de Dados, Machine Learning, Python.
\end{resumo}

\newpage

%%%%%%%%%%%%%%%%%%%%%%%%%%%%
%%                        %%
%%        Abstract        %%
%%                        %%
%%%%%%%%%%%%%%%%%%%%%%%%%%%%


\begin{abstract}
    \noindent
    \justifying

    This work aims to present a proposal for the development of an extensible and customizable data analysis and visualization platform, aiming to provide a simple and objective way of accessing machine learning libraries and the like present in the library repository \textit{Pypi }. To achieve this goal, interviews were conducted with researchers in the area of Data Science and Signal Processing, responsible for the research project \textit{SmartEnergy}, to identify the main needs and challenges in the data analysis process. In addition, a literature search was carried out to identify the main data analysis and data visualization tools that can meet the identified requirements. Finally, a set of libraries and tools for the development of the platform was proposed, which can be used to create a development environment for data analysis and visualization, which meets the needs of the data analysis process.


    \textbf{Keywords}: Platform, Data Analysis, Data Visualization, Machine Learning, Python.
\end{abstract}

\newpage
